\documentclass[a4paper,12pt]{article}
\usepackage[UTF8,fontset=none]{ctex}
\usepackage[top=2.5cm,bottom=2cm,left=2.5cm,right=2cm]{geometry}
\usepackage{array}
\usepackage{fontspec}

\setmainfont{TeX Gyre Termes}
\setCJKmainfont{Source Han Serif CN}
\setCJKsansfont{Noto Sans CJK SC}

\pagestyle{empty}
\setlength{\parindent}{0pt}
\renewcommand{\arraystretch}{2.0}

\newlength{\TW}
\setlength{\TW}{\dimexpr\textwidth-2\tabcolsep-2\arrayrulewidth}

\begin{document}

\begin{center}
{\Large\bfseries 湖北文理学院毕业论文（设计）档案材料清单}
\end{center}

\vspace{10pt}

\noindent
\begin{tabular}{|p{2cm}|p{4cm}|p{2cm}|p{\dimexpr\TW-8cm-6\tabcolsep-2\arrayrulewidth}|}
\hline
题目 & \multicolumn{3}{p{\dimexpr\TW-2cm-2\tabcolsep}|}{基于Apollo的多障碍物场景下路径规划算法实现} \\
\hline
专业班级 & \multicolumn{3}{p{\dimexpr\TW-2cm-2\tabcolsep}|}{计科2211} \\
\hline
学生姓名 & 廖嘉旺 & 学\hspace{1em}号 & 2022117117 \\
\hline
指导教师 & 孙成娇 & 职\hspace{1em}称 & 副教授 \\
\hline
\end{tabular}

\vspace{16pt}

\noindent
\begin{tabular}{|p{2cm}|p{\dimexpr\TW-2cm-2\tabcolsep}|}
\hline
序号 & \centering 材料名称 \tabularnewline
\hline
1 & 学生论文 \\
\hline
2 & 任务书 \\
\hline
3 & 开题报告 \\
\hline
4 & 学生论文查重报告 \\
\hline
5 & 学生论文规范性检测报告 \\
\hline
6 & 答辩资格审查表 \\
\hline
7 & 教师指导记录 \\
\hline
8 & 学生论文答辩情况记录表 \\
\hline
9 & 学生毕业论文（设计）成绩评定表 \\
\hline
\end{tabular}

\end{document}
